\section{Implementation and evaluation details}\label{app:protocol}
\textbf{Supervision implementation.} A drug record is a compound entry within a source and cell line. MOA masks are constructed within each sampled batch; unannotated drugs can still contribute to compound pairing. For cpg0000 compounds, MOA is assigned as the most frequent reference MOA among drugs targeting the associated gene. $L_1,L_2,L_3,L_4$ have separate learned temperatures; $L_5$ shares the $L_2$ temperature. Only these five objectives are active. Within-modality losses average over anchors with positives; cross-group losses average over all anchors, with zero contribution from anchors without positives.

\textbf{Retrieval settings.} Same-compound retrieval uses the compound head, with five random 20\% query/80\% reference crop splits. Functional retrieval uses MOA heads and excludes the query drug in within-cell-line and cross-batch MOA searches. U2OS and A549 query each other's reference libraries. Cross-batch retrieval uses 115 cpg0030 queries and 690 cpg0012 references; cpg0030 queries are excluded from encoder training. Batch mixing uses 32 shared compounds. HT29 images are excluded from training; 71 drug queries search 1,729 pooled U2OS/A549 records. CRISPR retrieval uses 29 of the 40 held-out A549 genes with positives among 1,038 reference drugs. Positives are annotated target drugs plus drugs sharing an MOA with those drugs, expanded once. Functional queries must have at least one positive.

\textbf{Metrics and statistics.} Random rates use query-specific library sizes and positive counts. Batch-mixing scores are ASW, iLISI, Leiden-cluster entropy and normalized Gini, oriented so that higher is better; panel-specific settings are stated in captions.

Query-bootstrap 95\% CIs use the 2.5th/97.5th percentiles, averaging repeated hit indicators within each query first; models and libraries are fixed. Panels showing standard deviations are identified in captions. Figs.~2B, 3A and 4 use paired-query two-sided Wilcoxon signed-rank tests: zero differences are discarded, tied absolute differences receive average ranks, and exact conditional sign-randomization $p$-values are Holm-adjusted within families of 26, 4 and 52 comparisons, respectively (the last includes both heads). In Fig.~3A, the five crop draws are averaged within each checkpoint and then equally across the three checkpoints for each of the 115 queries; the four tests compare the MOA head with the four displayed baselines. Pratt and paired sign-flip tests provide sensitivity checks.

\textbf{Parameter counting.} Table~\ref{tab:model-size} counts parameters in the image inference path, including frozen weights, with M denoting $10^6$ parameters. BioMoAX includes the shared ViT and one active projection head; its compound and MOA heads have equal sizes. DeepProfiler excludes its unused training classifier. CLOOME includes the image output layer; CellCLIP counts the DINOv2-G backbone once across channels, together with the channel Transformer. Expression, chemical and text encoders, unused heads, non-parameter buffers and ridge mappings are excluded. CellProfiler has no neural image encoder, and raw L1000 expression profiles are not part of this image-side comparison.

\textbf{Representation visualization.} Fig.~2C highlights the same annotated crops across methods, sampling 20 MOAs, four drugs per MOA and three crops per drug. In Fig.~4C, UMAP is fitted on A549 and applied to U2OS and HT29. Seven distant background points are omitted from each BioMoAX and CellProfiler map; all highlighted drug profiles are retained.

\textbf{Case visualization.} Networks in Fig.~5B and C use seed45\_ep028 and seed43\_ep029, respectively. Each network includes all query genes and the union of their Top-10 reference drugs across five crop draws for that checkpoint. Both ALDH2 and AKR1B1 retrieve disulfiram in all five crop draws of seed43, but not among the Top-10 in the other two checkpoints. Disulfiram is also retrieved broadly by other genes in seed43. The examples therefore illustrate shared-drug retrieval, not gene specificity or new direct-target validation. The pink edges indicate the current reference annotations, whose absence for the dashed edges does not establish absence from all literature. External evidence for DDR2--linifanib comes from an in vitro kinase-binding screen~\cite{karaman2008}. The AKR1B1 protein decrease reported by Yoshino et al. was observed in renal carcinoma xenografts and was not reproduced in the same study\textquotesingle s cultured-cell experiments~\cite{yoshino2020npl4}.
